% Source printed pp. 141-144, physical PDF pp. 149-152.
\section{Deligne--Mumford 叠的刻画}\label{src:4.6-section}
在证明稳定子光滑的有限型点附近存在\newterm{极小呈示}{minimal presentation} 之后, 我们证明 Deligne--Mumford 叠与代数空间的刻画, 并由此推出 $\mathcal M_g$ 是 Deligne--Mumford 叠.
\subsection{极小呈示的存在性}\label{src:4.6.1-subsection}
\begin{theorem}[极小呈示的存在性]\label{src:4.6.1}
设 $\mathcal X$ 是 Noether 代数叠, $x\in|\mathcal X|$ 是稳定子 $G_x$ 光滑的有限型点. 则存在概形 $U$, 闭点 $u\in U$, 以及从概形 $U$ 出发, 相对维数为 $\dim G_x$ 的光滑态射 $(U,u)\to(\mathcal X,x)$, 使下图为 Cartesian 图:
\[
\begin{tikzcd}
\Spec\kappa(u)\arrow[r,hook]\arrow[d]&U\arrow[d]\\
\mathcal G_x\arrow[r,hook]\arrow[ur,phantom,"\square" description]&\mathcal X.
\end{tikzcd}
\]
特别地, 若 $G_x$ 有限且约化, 则存在从概形出发的平展态射 $(U,u)\to(\mathcal X,x)$.
\end{theorem}
\begin{proof}
设 $(U,u)\to(\mathcal X,x)$ 是从概形 $U$ 出发, 相对维数为 $n$ 的光滑态射, 且 $u\in U$ 是有限型点. 由命题~\bookref{4.5.16}, $x$ 处的剩余 gerbe $\mathcal G_x$ 存在, 包含映射 $\mathcal G_x\hookrightarrow\mathcal X$ 是局部闭浸入, 且 $\mathcal G_x$ 正则, 维数为 $-\dim G_x$ (习题~\bookref{4.5.22}). 得到 Cartesian 图
\[
\begin{tikzcd}
O(u)\arrow[r,hook]\arrow[d]&U\arrow[d]\\
\mathcal G_x\arrow[r,hook]\arrow[ur,phantom,"\square" description]&\mathcal X,
\end{tikzcd}
\]
其中 $O(u)$ 是维数为 $c:=n-\dim G_x$ 的正则概形. 设 $f_1,\ldots,f_c\in\mathcal O_{O(u),u}$ 是生成 $u$ 处极大理想的正则序列. 将 $U$ 替换为 $u$ 的一个仿射开邻域后, 可以假设每个 $f_i$ 都是 $U$ 上的整体函数. 闭子概形 $W:=V(f_1,\ldots,f_c)$ 在 $U$ 处与 $O(u)$ 横截相交, 即在概形意义下有 $W\cap O(u)=\Spec\kappa(u)$. 我们断言, 归纳地应用平坦性的切片判据~\bookeqref{A.2.9}, 可推出复合 $W\hookrightarrow U\to\mathcal X$ 在 $u$ 处平坦. 事实上, 由光滑下降, 只需验证 $W\times_{\mathcal X}U\hookrightarrow U\times_{\mathcal X}U\xrightarrow{p_2}U$ 在 $u$ 关于 $p_1:W\times_{\mathcal X}U\to W$ 的一个原像 $u'$ 处平坦. 由于诱导映射 $\Spec\kappa(u)\to\mathcal G_x$ 平坦, 其基变换 $U\times_{\mathcal X}\Spec\kappa(u)\to O(u)$ 也平坦, 因而 $f_1,\ldots,f_c$ 的拉回在 $U\times_{\mathcal X}\Spec\kappa(u)$ 于 $u'$ 处的局部环中构成正则序列. 归纳地应用平坦性的切片判据便得到这一断言.

由于 $G_x$ 光滑, $\Spec\kappa(u)\to\mathcal G_x$ 也光滑. 对平坦态射而言, 光滑性可以在纤维上检验, 因而 (再次通过光滑下降论证) $W\to\mathcal X$ 在 $u$ 处光滑. 将 $W$ 替换为 $u$ 的一个开邻域后, 即得结论.
\end{proof}
\begin{remark}\label{src:4.6.2}
若光滑呈示 $p:U\to\mathcal X$ 满足 $T_{U,u}\to T_{\mathcal X,p(u)}$ 是 $k$-向量空间同构, 则称它在 $u\in U(k)$ 处\newterm{极小万有}{miniversal}. 我们将在命题~\bookref{4.7.6} 中看到, 上面的呈示都是极小万有的.

若稳定子 $G_x$ 不光滑, 则有两种候选的 ``极小呈示''. 仍然存在极小万有呈示 $(U,u)\to(\mathcal X,x)$, 但其相对维数等于 $G_x$ 的 Lie 代数的维数 (而不是 $\dim G_x$), 因而纤维积 $\mathcal G_x\times_{\mathcal X}U$ 可能具有正维数. 例如, 在特征 $p$ 下, $B\mu_p$ 是代数叠 (命题~\bookref{7.3.10}), 并且可以实现为 $\mathbb G_m$ 通过 $t\cdot x=t^p x$ 作用于自身的商; 这里 $\mathbb G_m\to B\mu_p$ 是极小万有呈示. 另一方面, 上述论证可以推广, 证明存在 fppf 的 (但未必光滑的) 态射 $(U,u)\to(\mathcal X,x)$, 使 $\mathcal G_x\times_{\mathcal X}U\cong\Spec\kappa(u)$. 若进一步假设 $\mathcal X$ 具有拟有限对角线, 则 $U\to\mathcal X$ 是 fppf 且拟有限的呈示. 另见 \cite[Tag 06MC]{SP}.
\end{remark}
\begin{exercise}[技术性]\label{src:4.6.3}
若 $\mathcal X$ 是代数叠 (可以不是 Noether 的), $x\in\mathcal X$ 是稳定子 $G_x$ 离散且非分歧的有限型点, 证明存在从概形 $U$ 出发的平展态射 $(U,u)\to(\mathcal X,x)$, 其中 $u\in U$ 是闭点.

\textit{提示:} 利用习题~\bookref{4.5.18}, 重复上述论证.
\end{exercise}
\subsection{等价刻画}\label{src:4.6.2-subsection}
\begin{theorem}[Deligne--Mumford 叠的刻画]\label{src:4.6.4}
设 $\mathcal X$ 是代数叠. 下列条件等价:
\begin{enumerate}[(1)]
\item 叠 $\mathcal X$ 是 Deligne--Mumford 叠;
\item 对角线 $\mathcal X\to\mathcal X\times\mathcal X$ 非分歧;
\item $\mathcal X$ 的每个点的稳定子群都离散且约化.
\end{enumerate}
若 $\mathcal X$ 具有拟紧对角线 (例如 $\mathcal X$ 拟分离), 则 (3) 等价于要求每个稳定子都有限且约化.
\end{theorem}
\begin{proof}
等价关系 $(2)\Leftrightarrow(3)$ 本质上就是非分歧的定义: 由于对角线 $\mathcal X\to\mathcal X\times\mathcal X$ 总是局部有限型的 (习题~\bookref{4.3.5}), 它非分歧当且仅当每个几何纤维 (或为空, 或同构于一个稳定子) 都离散且约化. 不难证明 Deligne--Mumford 叠的对角线非分歧 (习题~\bookref{4.2.8}). 反过来, 极小呈示的存在性 (定理~\bookref{4.6.1} 和习题~\bookref{4.6.3}) 表明, 对每个有限型点 $x\in\mathcal X$, 都存在从概形出发的平展态射 $U\to\mathcal X$, 其像包含 $x$. 因此 $\mathcal X$ 是 Deligne--Mumford 叠. 另见 \cite[定理 $8.1$]{LMB00} 和 \cite[Tag 06N3]{SP}.
\end{proof}
\begin{corollary}\label{src:4.6.5}
代数叠的态射 $\mathcal X\to\mathcal Y$ 是相对 Deligne--Mumford 的 (定义~\bookref{4.3.36}), 当且仅当 $\mathcal X\to\mathcal X\times_{\mathcal Y}\mathcal X$ 非分歧.
\end{corollary}
\begin{theorem}[代数空间的刻画]\label{src:4.6.6}
设 $\mathcal X$ 是对角线可由概形表出的代数叠. 下列条件等价:
\begin{enumerate}[(1)]
\item 叠 $\mathcal X$ 是代数空间;
\item 对角线 $\mathcal X\to\mathcal X\times\mathcal X$ 是单态射;
\item $\mathcal X$ 的每个点的稳定子都平凡.
\end{enumerate}
\end{theorem}
\begin{remark}\label{src:4.6.7}
我们将在定理~\bookref{5.5.10} 中去掉 $\Delta_{\mathcal X}$ 可由概形表出这一恼人的假设.
\end{remark}
\begin{proof}
条件 (2) 等价于 $\mathcal X$ 是层 (习题~\bookref{3.4.39}(d)). 蕴涵 $(1)\Rightarrow(2)$ 由代数空间的定义得到. 反过来, 若 $\mathcal X$ 是层, 则极小呈示的存在性~\bookeqref{4.6.1} 蕴含存在从概形出发的满射, 平展且可表的态射 $U\to\mathcal X$. 由于 $\Delta_{\mathcal X}$ 可由概形表出, $U\to\mathcal X$ 也可由概形表出. 等价关系 $(2)\Leftrightarrow(3)$ 由以下事实得到: 有限型群概形平凡当且仅当每个纤维都平凡 (命题~\bookref{B.1.9}).
\end{proof}
\begin{corollary}\label{src:4.6.8}
设 $G\to S$ 是作用于 $S$ 上代数空间 $U$ 的光滑仿射群概形. 则
\begin{enumerate}[(1)]
\item $[U/G]$ 是 Deligne--Mumford 叠当且仅当 $U$ 的每个点的稳定子群都离散且约化, 或等价地, 当且仅当作用映射 $G\times U\to U\times U$ 非分歧.
\item $[U/G]$ 是代数空间当且仅当 $U$ 的每个点的稳定子群都平凡, 或等价地, 当且仅当作用映射 $G\times U\to U\times U$ 是单态射.
\end{enumerate}
\end{corollary}
\begin{corollary}[可表态射的刻画]\label{src:4.6.9}
设 $\mathcal X\to\mathcal Y$ 是对角线可由概形表出的 Noether 代数叠的态射. 则 $\mathcal X\to\mathcal Y$ 可表当且仅当对每个几何点 $x\in\mathcal X(k)$, 自同构群之间的映射 $G_x\to G_{f(x)}$ 都是单射.
\end{corollary}
\begin{corollary}\label{src:4.6.10}
对于 $g\geq2$, $\mathcal M_g$ 是 $\mathbb Z$ 上有限型且具有仿射对角线的 Deligne--Mumford 叠.
\end{corollary}
\begin{proof}
只需证明 $\mathcal M_g$ 是 Deligne--Mumford 叠. 由 Deligne--Mumford 叠的刻画~\bookeqref{4.6.4}, 只需证明域 $k$ 上每条光滑连通射影曲线 $C$ 的自同构群概形 $\Aut(C)$ 都离散且约化. 换言之, 需证明 Lie 代数 $T_{\Aut(C),e}$ 的维数为零. 考虑下列图:
\[
\begin{tikzcd}[column sep=3.4em]
\Spec k\arrow[r,"e"]\arrow[d,hook]&\Aut(C)\arrow[r]\arrow[d]&\mathcal M_g\arrow[d]\\
\Spec k[\epsilon]\arrow[r]\arrow[ur,"\tau"]&\Spec k\arrow[r,"{([C],[C])}"]\arrow[ur,phantom,"\square" description]&\mathcal M_g\times\mathcal M_g
\end{tikzcd}
\qquad\text{与}\qquad
\begin{tikzcd}[column sep=2.6em,row sep=3em]
\Spec k\arrow[r,"{[C]}"]\arrow[d,hook]&\mathcal M_g\arrow[d]\\
\Spec k[\epsilon]\arrow[r,"p"']\arrow[ur,""{name=a}]\arrow[ur,bend right=30,""{name=b}]&\Spec k
\arrow[from=a,to=b,Rightarrow,"\tau"']
\end{tikzcd}
\]
左图的提升 $\tau:\Spec k[\epsilon]\to\Aut(C)$, 即切向量 $\tau\in T_{\Aut(C),e}$, 在右图中对应于平凡一阶形变 $[C]\circ p$ 的一个自同构 $\tau$. 由形变理论~\bookeqref{C.2.4}, 平凡一阶形变的自同构群可等同于 $H^0(C,T_C)$. 但这个向量空间为零, 因为 $T_C=\Omega_C^\vee$ 的次数为 $2-2g<0$.
\end{proof}
